\subsection{Atributos derivados}

Los atributos derivados que construimos son la vista de cada foto, si la foto es usable, el BCS de entrenamiento, el número de evaluaciones, la diferencia entre evaluadores, el pliegue y el número de fotos usables por vista de cada vaca.

\subsubsection*{Vista y foto usable}

Ni los nombres ni los metadatos de las fotos dicen la vista, así que la marcamos a mano con la página \nolinkurl{src/bigdata/herramientas/marcar_vistas.html} del repositorio. Con una tecla se marca 1 posterior, 2 lateral izquierda, 3 lateral derecha o 0 descartar (borrosa, cortada, tapada o en una vista que no sirve). Izquierda y derecha son del lado de la vaca y no de la foto. Cada marca guarda quién la hizo y cuándo. Las fotos descartadas no se borran y quedan en una carpeta aparte. Las que no se marquen quedarían en \nolinkurl{sin_marcar} y no se usarían, aunque en nuestra corrida no quedó ninguna sin marcar.

\begin{table}[H]
  \centering
  \caption{Fotos por vista.}
  \label{tab:vistas}
  \begin{tabular}{lrr}
    \toprule
    Vista & Fotos & Vacas con al menos una foto usable \\
    \midrule
    Posterior & 257 & 68 \\
    Lateral izquierda & 146 & 40 \\
    Lateral derecha & 148 & 36 \\
    Descartadas & 29 & No aplica \\
    \midrule
    Total & 580 & 76 vacas con fotos usables \\
    \bottomrule
  \end{tabular}
\end{table}

\begin{itemize}
  \item Las fotos de la primera visita traen sobre todo laterales. De sus 247 fotos, 43 son posteriores, 88 laterales izquierdas, 95 laterales derechas y 21 descartadas, y con ellas las laterales pasan de 111 a 294.
  \item De las 76 vacas con fotos usables, 18 no tienen ninguna foto lateral y 8 no tienen foto posterior. Además, 18 tienen las dos laterales y 17 tienen las tres vistas.
  \item La mediana es de 6 fotos usables por vaca, con un mínimo de 1 y un máximo de 29.
  \item Las 29 fotos descartadas son de 14 vacas, que son la 1507, 1516, 1517, 1536, 1554, 1561, 1564, 1569 (6 fotos), 1578 (2), 2114, 6137, 6226 (2), 8709 (2) y 91579 (8). La 1554 tenía una sola foto y por eso no tiene fotos usables.
\end{itemize}

\subsubsection*{Etiqueta de BCS}

A algunas vacas las evaluaron dos o tres veces, ya sea expertos distintos el mismo día o en las dos visitas. Como acordamos en la tabla de riesgos de 2.4, decidimos no elegir una evaluación sobre otra y armamos las siguientes columnas.

\begin{itemize}
  \item \texttt{bcs\_1}, \texttt{bcs\_2} y \texttt{bcs\_3} guardan las evaluaciones en orden, primero las de la bitácora del 14 de octubre y después las de la hoja. El orden no indica qué experto dio cada una.
  \item \texttt{n\_eval} es el número de evaluaciones de la vaca.
  \item \texttt{bcs\_entrenamiento} es el promedio de todas las evaluaciones de la vaca en las dos visitas y es la etiqueta con la que se entrena. Si solo hay una evaluación, es esa. Todas las fotos de la vaca usan este valor, sin importar la fecha de la foto.
  \item \texttt{dif\_evaluadores} es la diferencia entre la evaluación mayor y la menor, cuando hay dos o más.
\end{itemize}

Para evaluar el modelo se compara contra cada evaluación real (\nolinkurl{evaluaciones.csv}) y no contra el promedio. Las 27 vacas con dos o más evaluaciones que siguen en el dataset tienen una diferencia media de 0.185 y máxima de 0.50, y en 11 de ellas todas las evaluaciones son iguales.

\begin{table}[H]
  \centering
  \caption{Vacas por BCS de entrenamiento.}
  \label{tab:bcs-entrenamiento}
  \begin{tabular}{lr}
    \toprule
    BCS de entrenamiento & Vacas \\
    \midrule
    2.000 & 1 \\
    2.250 & 2 \\
    2.500 & 12 \\
    2.625 & 3 \\
    2.750 & 14 \\
    3.000 & 13 \\
    3.125 & 5 \\
    3.250 & 8 \\
    3.375 & 4 \\
    3.500 & 8 \\
    3.750 & 6 \\
    \midrule
    Total & 76 \\
    \bottomrule
  \end{tabular}
\end{table}

Los valores como 2.625, 3.125 y 3.375 salen de promediar evaluaciones distintas, y el 2.0 es de una sola vaca (1213).

\subsubsection*{Pliegues de validación cruzada}

Repartimos las vacas en 5 pliegues ($k = 5$). En CRISP-DM esto es la tarea 4.2 (diseño de prueba), pero lo hacemos aquí para que el dataset salga listo para entrenar.

\begin{itemize}
  \item La separación es por vaca y no por foto, así que todas las fotos de una vaca quedan en el mismo pliegue. Si una vaca tuviera fotos en entrenamiento y en prueba, el modelo podría aprender a reconocerla en lugar de aprender su BCS y las métricas saldrían infladas (\textit{data leakage}).
  \item Los pliegues no son aleatorios. Ordenamos las vacas por BCS de entrenamiento (y por ID cuando empatan) y las asignamos en rotación (0, 1, 2, 3, 4, 0, 1 y así), de modo que cada pliegue tiene BCS bajos, medios y altos y el reparto sale igual en cada corrida.
  \item Una vaca recibe pliegue solo si tiene al menos una foto usable y BCS.
\end{itemize}

\begin{table}[H]
  \centering
  \caption{Vacas, fotos y BCS medio por pliegue.}
  \label{tab:pliegues}
  \begin{tabular}{lrrrrrr}
    \toprule
    Pliegue & Vacas & Fotos & Posterior & Lateral izq. & Lateral der. & BCS medio \\
    \midrule
    0 & 16 & 146 & 67 & 43 & 36 & 2.99 \\
    1 & 15 & 69 & 42 & 16 & 11 & 2.96 \\
    2 & 15 & 113 & 49 & 25 & 39 & 2.98 \\
    3 & 15 & 138 & 61 & 39 & 38 & 3.02 \\
    4 & 15 & 85 & 38 & 23 & 24 & 3.02 \\
    \bottomrule
  \end{tabular}
\end{table}

Los pliegues quedan parejos en vacas y en BCS, pero no en número de fotos, ya que el pliegue 1 tiene 69 y el 0 tiene 146. Esto pasa porque algunas vacas tienen muchas fotos, por ejemplo la 1569 tiene 29 usables y la 1561 tiene 26. Como la separación es por vaca, en el modelado conviene revisar que esas vacas no pesen de más en el entrenamiento y en la evaluación.

\subsubsection*{Otras decisiones}

\begin{itemize}
  \item En esta fase no normalizamos, es decir, no cambiamos la resolución ni los valores de los píxeles. Las fotos se quedan a resolución completa y se reducen después, al entrenar.
  \item Las 73 fotos del dataset final cuyo HEIC traía mapa de profundidad del LiDAR solo se registran en la columna \texttt{profundidad}, y los mapas no se guardan.
\end{itemize}
